% Source: https://artofproblemsolving.com/community/c3h2128640_2018_mathcounts_national_number_theory_problem
% Stable identity: aops-topic-2128640-q01
Since $2018\equiv18\equiv-2\pmod{20}$, we need $(-2)^{2018}=2^{2018}\pmod{20}$. Work modulo $4$ and $5$:
\begin{itemize}
\item modulo $4$: $2^{2018}\equiv0$, since $2018\ge2$;
\item modulo $5$: $2^4=16\equiv1$, and $2018=4\cdot504+2$, so $2^{2018}\equiv2^2=4$.
\end{itemize}
The number $r$ with $0\le r<20$, $r\equiv0\pmod4$ and $r\equiv4\pmod5$ is $r=4$ (the candidates $0,4,8,12,16$ have remainders $0,4,3,2,1$ modulo $5$). So the remainder is $\boxed{4}$.

(Alternatively, the powers $2^2,2^3,2^4,2^5,2^6,\dots$ have remainders $4,8,16,12,4,\dots$ modulo $20$, repeating with period $4$; since $2018-2$ is a multiple of $4$, $2^{2018}\equiv2^2=4$.)
